\scanpage{8}
\deleted{revient à la donnée d’un scindage}
\diagram{\deleted{\begin{tikzcd}[ampersand replacement=\&]
\unsure{\mathcal G|\mathscr G}'\arrow[r]\&\Gamma\arrow[l,bend left=45]
\end{tikzcd}}}

\textcircled{2.3}. Soit
\[
\unsure{\mathfrak G|G|\mathscr G}'=\mathrm{Aut}_{\bar k}(\unsure{\mathscr X|\mathcal X|\mathfrak X}'\longrightarrow\mathscr X)\tag*{(17)}
\]
\deleted{\unsure{Alors} $\mathfrak G'$ se \unsure{dé|d’}} \unsure{donc|une|dans} suite exacte
\diagram{(18)\quad\begin{tikzcd}[ampersand replacement=\&,column sep=small]
1\arrow[r]\&G\arrow[r,dash]\&\mathfrak G'\arrow[r]\&\mathfrak G.
\end{tikzcd}\quad,}
qui se déduit (\unsure{mod.|à un|via l’} iso. can) de (16) comme
\[
\begin{aligned}
\mathfrak G'&\simeq\text{image inverse de }\mathfrak G\subset\mathcal G\\
&\hspace{1em}\text{par }\mathcal G'\mathbin{-}\mathcal G\\
&\simeq\mathrm{Ker}(\mathcal G'\mathbin{-}\Gamma)
\end{aligned}\tag*{(19)}
\]

Le cas favorable $N$ inv. dans $\unsure{\mathcal E|\mathscr E}$ i.e. $\mathcal G'\mathbin{-}\mathcal G$ surjectif revient à dire que a) $\mathfrak G'\mathbin{-}\mathfrak G$ surjectif et b) $\mathcal G'\mathbin{-}\Gamma$ \deleted{surjectif} $\overline{\mathrm{m}}$ image que $\mathcal G\to\Gamma$

\deleted{\uline{\unsure{NB}}}

\uline{NB} Si $\mathscr X$ de prés. finie sur $\unsure{S|\mathcal S}$, alors l’image $\Gamma_1$ de $\deleted{\unsure{\Gamma}}\mathcal G$ dans $\Gamma$ est un sous-groupe ouvert, et si de plus $\mathscr X'$ est fini sur $\mathscr X$ iee de prés. finie sur $k$, alors l’image $\Gamma'_1$ \inserted{$(\subset\Gamma_1)$} de $\mathcal G'$ dans $\Gamma$ l’est également.
% archive sheet number: 7
